\documentclass[10pt,a4paper]{amsart}
\usepackage[margin=19mm]{geometry}
\usepackage{amsmath,amssymb,mathtools}
\usepackage[scr=rsfs,cal=euler]{mathalfa}
\usepackage{xfrac}
\usepackage[unicode,hidelinks,bookmarks=false]{hyperref}
\newcommand{\ie}{i.e.\ }
\newcommand{\cf}{cf.\ }
\newcommand{\resp}{respectively\ }
\newcommand{\Subsection}{\subsection}
\providecommand{\nobreakdash}{\nobreak\hbox{-}}
\setlength{\emergencystretch}{2em}
\multlinegap0pt
\usepackage{bm}
\usepackage{bbm}
\usepackage{dsfont}
\usepackage{xspace}
\usepackage{cancel}
\usepackage{mathtools}
\usepackage{amsthm}
\makeatletter
\@ifundefined{theorem}{\newtheorem{theorem}{Theorem}}{}
\@ifundefined{proposition}{\newtheorem{proposition}[theorem]{Proposition}}{}
\makeatother
\newcommand{\Chain}[1]{\mathscr{C}^{#1}}
\newcommand{\F}{{\rm F}}
\newcommand{\HH}[2]{H^{1,#1}_{{\rm #2}}}
\newcommand{\LUG}[1]{{\rm LUG}^{#1}}
\newcommand{\Lip}{{\rm Lip}}
\newcommand{\N}{\mathbb{N}}
\newcommand{\R}{\mathbb{R}}
\newcommand{\UG}[1]{{\rm UG}^{#1}}
\newcommand{\X}{{\rm X}}
\newcommand{\limi}{\varliminf}
\newcommand{\lims}{\varlimsup}
\newcommand{\lip}{{\rm lip}}
\newcommand{\mm}{\mathfrak m}  
\newcommand{\restr}[1]{\lower3pt\hbox{$|_{#1}$}}
\newcommand{\sfc}{{\sf c}}
\newcommand{\sfd}{{\sf d}}
\renewcommand{\d}{{\mathrm d}}
\title[Sobolev spaces via chains in metric measure spaces]{Sobolev spaces via chains in metric measure spaces}
\author{Emanuele Caputo, Nicola Cavallucci}
\date{Source: arXiv:2408.15071v1 (CC BY 4.0)}
\begin{document}
\maketitle

\noindent\emph{Source and licence.} Sobolev spaces via chains in metric measure spaces. Version arXiv:2408.15071v1 (CC BY 4.0), distributed under the Creative Commons Attribution 4.0 International licence, as recorded in the arXiv licence metadata for this exact version and shown on the abstract page https://arxiv.org/abs/2408.15071v1. Full rights notice: licenses/arxiv-2408.15071-CC-BY-4.0.txt.\par\smallskip

\noindent\textit{Editorial scope.} Counted unit: the Theorem whose source label is \texttt{theo:density-energy-chains-complete} in the article (source0.tex lines 1493--1504, source label \texttt{theo:density-energy-chains-complete}), delivered together with the article's own complete original proof of it (source0.tex lines 1549--1601, one \texttt{proof} environment of 53 lines). No mathematics is written, repaired or re-derived: statement and proof are the article's own text line for line. The proof is detached from the statement in the source: the article states the result at lines 1493--1504 and prints its proof at lines 1549--1601, and the proof environment's own header names the statement, so the association is the article's own. Required context retained uncounted: prop:UGeta\_subset\_UG (lines 761--769); theo:density\_energy\_complete (lines 1241--1251); prop:Leibniz-rule-chains (lines 1507--1522). Every definition, display and label of those lines is retained unchanged. Omitted from this package and not redistributed: 1--760, 770--1240, 1252--1492, 1505--1506, 1523--1548, 1602--2402. Nothing inside a retained excerpt is omitted or altered: every retained body is delivered exactly as the article's lines are written, so no omission receipt of any kind accompanies this package. Citations: the retained excerpts carry no citation command at all, so no bibliography entry of the article is transcribed and no reference is redistributed. Reference adaptation: none was needed and none was made; every reference inside the delivered unit resolves inside it, so no unresolved reference or citation is produced. Printed slips kept literal: no printed word, symbol, hypothesis or number of the counted statement or of its proof was altered. Layout and class: the article's own class is \texttt{amsart} with packages including geometry, inputenc, babel, fontenc, mathrsfs, mathtools, amsmath, amssymb, epsfig, caption, subcaption, amsthm, textpos, emptypage; the wrapper is the lane's pinned \texttt{amsart} editorial wrapper at 10pt on A4, which restates the theorem-like environments the retained text needs and adds the article's own macro definitions that the retained text needs (\texttt{\textbackslash Chain}, \texttt{\textbackslash F}, \texttt{\textbackslash HH}, \texttt{\textbackslash LUG}, \texttt{\textbackslash Lip}, \texttt{\textbackslash N}, \texttt{\textbackslash R}, \texttt{\textbackslash UG}, \texttt{\textbackslash X}, \texttt{\textbackslash d}, \texttt{\textbackslash limi}, \texttt{\textbackslash lims}, \texttt{\textbackslash lip}, \texttt{\textbackslash mm}, \texttt{\textbackslash restr}, \texttt{\textbackslash sfc}, \texttt{\textbackslash sfd}), with extra wrapper packages (bm, bbm, dsfont, xspace, cancel, mathtools). The delivered numbering is the unit's own. Assets: no retained span contains a figure environment, a graphics-inclusion command, a PDF-page inclusion, a table or a TikZ picture; no image file is copied into this package, the package declares no asset dependency and no image file is redistributed with it. Licence: version arXiv:2408.15071v1 of this article is distributed under the Creative Commons Attribution 4.0 International licence, as recorded in the arXiv licence metadata for this exact version and shown on the abstract page https://arxiv.org/abs/2408.15071v1. A full rights notice is delivered at licenses/arxiv-2408.15071-CC-BY-4.0.txt. Characters: every retained excerpt is pure ASCII, so the delivered engine, XeLaTeX with its default Latin Modern text font, reports no missing character.
\begin{proposition}
\label{prop:UGeta_subset_UG}
    Let $(\X,\sfd)$ be a metric space and let $u\colon \X \to \R$. If $g\in \UG{\varepsilon}(u)$ for some $\varepsilon >0$, then 
    \begin{equation}
    \label{eq:UG_inequality}
        \vert u(\omega(\gamma)) - u(\alpha(\gamma)) \vert \le \int_\gamma g
    \end{equation}
    for every rectifiable curve $\gamma$ such that $\omega(\gamma), \alpha(\gamma) \in \{g < +\infty \}$. In particular if $g$ takes only finite values then $g \in \UG{}(u)$.
\end{proposition}

\begin{theorem}
    \label{theo:density_energy_complete}
    Let $(\X,\sfd,\mm)$ be a metric measure space such that $(\X,\sfd)$ is complete. Then \begin{equation}
        \HH{p}{\Chain{},\, \Lip}(\X) = \HH{p}{curve}(\X)
    \end{equation}
    and
    \begin{equation}
        \|u \|_{\HH{p}{\Chain{},\, \Lip}(\X)}  = \|u \|_{\HH{p}{curve}(\X)}
    \end{equation}
    for every $u\in L^p(\X)$.
\end{theorem}

\begin{proposition}[Leibniz rule]
\label{prop:Leibniz-rule-chains}
    Let $u\colon \X \to \R$ be Borel and $\varphi \in {\rm Lip}(\X)$.
    For every $g \in {\rm UG}^\varepsilon(u)$, we have
    \begin{equation}
    \label{eq:varepsilon_upper_gradient_Leibniz}
        \vert u \vert\, {\rm sl}_\varepsilon \varphi + Q_\varepsilon \varphi\, g \in {\rm UG}^\varepsilon(u \varphi),
    \end{equation}
    where $Q_\varepsilon\varphi(x):= \sup_{y\in \overline{B}_\varepsilon(x)} \vert \varphi \vert(y)$. In particular, for every $u \in \mathcal{L}^p(\X)$ it holds that
    \begin{equation}
    \label{eq:Leibniz_rule}
        \F{\Chain{}}(u\varphi) \le \left( \int |u|^p (\lip\,\varphi)^p \,\d \mm \right)^{\frac{1}{p}} + \|\varphi\|_{L^\infty(\X)}\, \F{\Chain{}}(u).
    \end{equation}

    
\end{proposition}

\begin{theorem}
    \label{theo:density-energy-chains-complete}
    Let $(\X,\sfd,\mm)$ be a metric measure space such that $(\X,\sfd)$ is complete. Then 
    \begin{equation}
        \HH{p}{\Chain{},\, \Lip}(\X) = \HH{p}{\Chain{}}(\X)
    \end{equation}
    and
    \begin{equation}
        \|u \|_{\HH{p}{\Chain{},\, \Lip}(\X)}  = \|u \|_{\HH{p}{\Chain{}}(\X)}
    \end{equation}
    for every $u\in L^p(\X)$.
\end{theorem}

\begin{proof}[Proof of Theorem \ref{theo:density-energy-chains-complete}]
    The proof is a modification of the one of Theorem \ref{theo:density_energy_complete}. We highlight what are the differences in each step.
    
    {\bf Step 1.} The proof does not change if we show that for every $u\in L^p(\X)$ we can find a sequence of bounded functions $u_j$ with bounded support such that $u_j \to u$ in $L^p(\X)$ and $\limi_{j\to +\infty}\F{\Chain{}}(u_j) \le \F{\Chain{}}(u)$. We fix a basepoint $x_0 \in \X$ and we consider the $1$-Lipschitz function $\varphi_j(x) = \max\{0,\min\{1,j+1-\sfd(x,x_0)\}\}$. We define $u_j := \min \{j, \max\{-j, \varphi_j u\}\}$. By definition, $u_j$ is bounded and has bounded support. Moreover, every $\varepsilon$-upper gradient of $\varphi_j u$ is also a $\varepsilon$-upper gradient of $u_j$. This implies that $\F{\Chain{}}(u_j) \le \F{\Chain{}}(\varphi_j u)$ for every $j$. Proposition \ref{prop:Leibniz-rule-chains} implies that
    \begin{equation}
        \begin{aligned}
            \limi_{j\to +\infty} \F{\Chain{}}(u_j) \le \limi_{j\to +\infty} \F{\Chain{}}(\varphi_j u) &\le \limi_{j\to +\infty}\left( \int |u|^p (\lip\,\varphi_j)^p \,\d \mm \right)^{\frac{1}{p}} + \|\varphi_j\|_{L^\infty}\, \F{\Chain{}}(u)\\
            &= \limi_{j\to +\infty}\left( \int_{\overline{B}_{j+1}(x_0)\setminus B_j(x_0)} |u|^p \,\d \mm \right)^{\frac{1}{p}} +\F{\Chain{}}(u) = \F{\Chain{}}(u),
        \end{aligned}
    \end{equation}
    where in the last equality we used that $u\in L^p(\X)$.
    
    {\bf Step 2.} It does not change. In particular the claim we have to prove is the following. For every $\varepsilon > 0$, for every $\varepsilon$-upper gradient $g\in \UG{\varepsilon}(u)$ and for every $\eta > 0$ there exists another $\varepsilon$-upper gradient $g_\eta \in \UG{\varepsilon}(u)$ such that 
    \begin{equation}
        \label{eq:approx_with_g_eta_chains}
        \Vert g - g_\eta \Vert_{L^p(\X)} < \eta
    \end{equation} 
    and with the following property. For every $j\in \N$ there exist functions $u_{\eta,j} \colon \X \to \R$ and $g_{\eta,j}\in \LUG{\frac{1}{j}}(u_{\eta,j})$ such that 
    \begin{equation}
        \label{eq:approx_with_u_eta,j_chain}
        \lims_{j\to +\infty} \Vert u_{\eta,j} - u \Vert_{L^p(\X)} \le 2\eta \quad \text{and} \quad \lim_{j\to +\infty} \Vert g_{\eta,j} - g_\eta \Vert_{L^p(\X)} = 0.
    \end{equation}
    Let $R\ge 1$ be such that $u\equiv 0$ on $\X\setminus B_{R}(x_0)$. We recall that it is enough to consider chain upper gradients $g\in \UG{\varepsilon}(u)$ that are lower semicontinuous and such that $g \equiv 0$ on $\X \setminus B_{2R}(x_0)$, by a truncation argument.

    {\bf Step 3.} The definition of $g_\eta$ does not change. Observe that $g_\eta \equiv 0$ on $\X \setminus B_{2R}(x_0)$ and that $g_\eta \in \UG{\varepsilon}(u)$ because $g_\eta \ge g$.
    
    {\bf Step 4.} The definition of $\hat{g}_{\eta,j}$ does not change and satisfies the same properties.
    
    {\bf Step 5.} Here there is a difference in the definition of the set $A$. Since $g_\eta \in L^p(\X)$ then $\mm(\{g_\eta = +\infty \})=0$. By outer regularity of the measure we can find an open set $U_\eta$ containing $\{g_\eta = +\infty \}$ and such that $\mm(U_\eta)^{\frac{1}{p}} < \eta (2M)^{-1}$. Moreover, since $g_\eta \equiv 0$ on $\X \setminus B_{2R}(x_0)$, we can choose $U_\eta$ such that $U_\eta \subseteq B_{2R}(x_0)$.
    Now we change the definition of the set $A$ by setting $A := (K_N \cup (\X \setminus (B_R(x_0))) \setminus U_\eta$. It is still closed and $u\restr{A}$ is still continuous.
    Now, the definition of $\hat{u}_{\eta,j}$ does not change, except for the fact that we use this set $A$. Properties (a)-(e) continue to hold.

    {\bf Step 6.} The definitions of $u_{\eta,j}$ and $g_{\eta,j}$ do not change.
    
    {\bf Step 7.} Here we claim that it is enough to show that $u_{\eta,j}(x)$ converges to $u(x)$ for every $x\in K_N\setminus U_\eta$ as $j \to +\infty$ and that $\hat{u}_{\eta,j}$ converges uniformly to $0$ on $\X\setminus B_{2R}(x_0)$. Indeed if this is true we have
    \begin{equation}
        \begin{aligned}
            &\lim_{j\to +\infty} \int_\X \Vert u - u_{\eta,j} \Vert_{L^p(\X)} \,\d\mm \\
            &\le\lim_{j\to +\infty}\left(\int_{K_N \setminus U_\eta} \vert u - u_{\eta,j} \vert^p \,\d\mm + \int_{(B_{2R}(x_0)\setminus K_N) \cup U_\eta} \vert u - u_{\eta,j} \vert^p \,\d\mm + \int_{(\X \setminus B_{2R}(x_0))} \vert u - u_{\eta,j} \vert^p \,\d\mm \right)^{\frac{1}{p}}\\
            &\le 0 + (2M) \mm(B_{2R}(x_0) \setminus K_N)^{\frac{1}{p}} + (2M)\mm(U_\eta)^{\frac{1}{p}} + 0 \le 2\eta.
        \end{aligned}
    \end{equation}

    {\bf Step 8.} Here we need an additional argument that justifies the different choice of $A$. Indeed, we claim that in any of three cases, the limit curve has its extreme points in $A$.
    
    In cases (1) and (3) this is true because either $\alpha(\sfc_j^e) = \alpha(\sfc_j) \in A$ by definition, or $\sfd(\alpha(\sfc_j^e), \X \setminus B_{2R}(x_0)) \le \frac{2}{j}$, by maximality of $\sfc_j^e$. In the first case the result is trivial because $A$ is closed and $u\restr{A}$ is continuous. In the second case $\alpha(\gamma) = \lim_{j\to +\infty} \alpha(\sfc_j^e) \in \X \setminus B_{2R}(x_0)$, so $\alpha(\gamma)\in A$ since $U_\eta \subseteq B_{2R}(x_0)$. Moreover, $u(\alpha(\gamma)) = 0 = u(\alpha(\sfc_j^e))$ because all these points belong to $\X \setminus B_R(x_0)$. On the other hand $\omega(\gamma) = \lim_{j\to +\infty} \omega(\sfc_j^e) = x \in A$ because $x$ is chosen in $K_N\setminus U_\eta$. Here, we have $u(\omega(\gamma)) = u(x) = u(\omega(\sfc_j^e))$.   
    
    In case (2) we have that $\alpha(\gamma) = \lim_{j\to +\infty} \alpha(\sfc_j^s) \in A$, because $\alpha(\sfc_j^s) \in A$ for every $j$ and $A$ is closed. Moreover, $u(\alpha(\gamma)) = \lim_{j\to +\infty} u(\alpha(\sfc_j^s))$ since $u\restr{A}$ is continuous. On the other hand, either $\omega(\sfc_j^s) = x$ for every $j$, so $\omega(\gamma) = x \in A$ and $u(\omega(\gamma)) = u(x) = u(\omega(\sfc_j^s))$, or $\sfd(\omega(\sfc_j^s), \X \setminus B_{2R}(x_0)) \le \frac{2}{j}$, by maximality of $\sfc_j^s$. Arguing as before, we get that $\omega(\gamma) \in \X \setminus B_{2R}(x_0)$, so it belongs to $A$ and $u(\omega(\gamma)) = 0 = u(\omega(\sfc_j^s))$. 
    
    In every case, the extreme points of $\gamma$ belong to the set $\{g_\eta < +\infty\}$. Hence $g_\eta$ satisfies the upper gradient inequality along $\gamma$ because of Proposition \ref{prop:UGeta_subset_UG}, while the proof shows that this is not the case, giving a contradiction.

    {\bf Step 9.} The proof does not change, using the same modifications we did in Step 8.
\end{proof}

\end{document}
